\section{Introduction}

3D Gaussian Splatting (3DGS) provides an explicit differentiable representation for dense scene rendering and has become a practical substrate for online visual world modeling~\cite{kerbl20233dgaussian}.
Gaussian SLAM systems such as Photo-SLAM, SplaTAM, GS-SLAM, and Gaussian-SLAM combine camera tracking with dense Gaussian reconstruction~\cite{huai2024photoSLAM,keetha2024splatam,yan2023gs,yugay2023gaussianslam}.
Their tracking frontends and persistent maps, however, still depend on a predominantly static world.
When people or objects occupy a large image region, dynamic observations can corrupt camera estimation and produce persistent Gaussian artifacts.

Semantic and geometric exclusion is therefore a natural defense.
Classical dynamic SLAM removes features associated with movable classes~\cite{bescos2018dynaslam,yu2018ds,xiao2019dynamic}, while recent dynamic Gaussian SLAM systems combine semantics, temporal consistency, depth warping, photometric evidence, or uncertainty to suppress unsafe observations~\cite{xu2024dgslam,kong2024dgsslam,liu2025dyphoslam,zhang2026dagsslam}.
This policy is appropriate for static map construction, but can be too restrictive for tracking.
A semantic mask can cover visible background, object boundaries, temporarily stationary surfaces, or pixels whose image motion is dominated by camera motion.
Removing all such observations can leave the pose estimator poorly constrained exactly when dynamic occlusion is strongest.

The key issue is not only whether an observation is static, but what authority it should receive.
Tracking and mapping have different failure costs and persistence.
A locally consistent observation may be useful for estimating the current camera pose while still being too uncertain to initialize or optimize persistent static structure.
Treating both decisions with one binary mask conflates \emph{tracking utility} with \emph{map eligibility}.

We propose \textbf{GDOR-SLAM}, a dynamic RGB-D Gaussian SLAM system based on Guarded Dynamic Observation Recovery (GDOR).
GDOR evaluates temporal depth consistency, residual optical-flow consistency, and the current tracking-risk state to recover a conservative subset of observations inside excluded dynamic regions.
Recovered observations may replenish features and support the current camera estimate.
They do not modify the conservative mask used for persistent ORB landmarks or Gaussian initialization, update, and densification.
The method therefore separates transient tracking privilege from persistent map privilege.

\begin{figure*}[t]
  \centering
  \includegraphics[width=\linewidth]{fig1_dyn19_overview.png}
  \caption{DYN-19 evidence overview. Full improves aggregate tracking over Semantic but not every sequence. The ordered configuration matrix identifies T+M as the strongest consumed-data tracking row. Multi-seed common-view reconstruction improves under Full, while the exact route audit distinguishes 30/30 non-vacuous zero-overlap Full passes from 30/30 expected Full-NoM overlap failures.}
  \label{fig:overview}
\end{figure*}

Our contributions are:
\begin{itemize}
    \item We instantiate an auditable two-channel observation contract for RGB-D/ORB Gaussian SLAM: tracking utility does not automatically grant persistent Gaussian map authority.
    \item We introduce GDOR, which combines temporal depth, residual-flow, and tracking-risk evidence to selectively recover dynamic-region observations for transient tracking.
    \item We bind every run to an exact recovered-support/map-admission overlap certificate and include a deliberate leakage counterfactual, making route separation empirically testable rather than only architectural.
    \item We report a complete 210-run ordered tracking matrix and a 72-cell multi-seed common-view mapping diagnostic, with sequence-level outcomes and proxy boundaries retained.
\end{itemize}

The supported claim is deliberately bounded.
The Full route improves aggregate tracking and reconstruction diagnostics over its same-source Semantic parent, and its route certificate passes in every tested cell with observed recovery.
It does not win every sequence, the mechanism matrix is ordered rather than factorial, and the leading T+M row was selected on consumed DYN-19 data.
The paper therefore argues for a concrete RGB-D/ORB recovery and routing mechanism rather than conceptual priority, universal improvement, external SOTA superiority, or direct measurement of 3D ghost contamination.
